\documentclass[a4paper]{article}
\usepackage[a4paper, top=8mm, bottom=8mm, left=8mm, right=8mm]{geometry}
\usepackage{polyglossia}
\setdefaultlanguage[babelshorthands=true]{russian}
\setotherlanguage{english}
\usepackage{fontspec}
\setmainfont{FreeSerif}
\newfontfamily{\russianfonttt}[Scale=0.7]{DejaVuSansMono}
\usepackage[tiny, compact]{titlesec}
\usepackage{titling}
\setlength{\droptitle}{-1cm}
\pretitle{\begin{center}\begin{bfseries}\Large}
\posttitle{\par\end{bfseries}\end{center}}
\preauthor{\begin{center}\normalsize}
\postauthor{\par\end{center}\vspace{-1.8cm}}
\usepackage{hyperref}
\usepackage{bookmark}
\usepackage{csquotes}
\usepackage{enumitem}
\usepackage{graphicx}
\usepackage{amsmath}

\title{Анализ графов(BFS): классическая реализация на C и алгоритм на основе GraphBLAS на примере parent BFS и multi source BFS}

\author{Пентегова Мария Сергеевна}

\date{}

\begin{document}

\begin{document}
\maketitle
\begin{flushright}
\begin{tabular}{r}
    Группа: \emph{\enquote{25.Б41-мм}} \\
    Кафедра: \emph{системного программирования} \\
    Научный руководитель: \emph{Семён Вячеславович} \\
    Номер семестра практики: \emph{2} \\
\end{tabular}
\end{flushright}

\section{Введение}

BFS (Breadth-First Search, обход в ширину)~--- один из фундаментальных алгоритмов теории графов. Он применяется для поиска кратчайших путей в невзвешенных графах, проверки связности и других задач.

Математически BFS от источника $s$ определяется следующим образом: множество вершин разбивается на уровни $L_0, L_1, L_2, \dots$, где:
\[
L_0 = \{s\}, \quad L_{k+1} = N(L_k) \setminus (L_0 \cup L_1 \cup \dots \cup L_k)
\]
где $N(L_k)$ --- множество всех соседей вершин из $L_k$.

В рамках учебной практики предлагается реализовать BFS двумя способами: классическим (на основе очереди) и с использованием библиотеки SuiteSparse:GraphBLAS, а затем сравнить их производительность.

\section{Постановка задачи}

\textbf{Целью работы является} реализация алгоритма BFS на языке C, а также его тестирование и сравнение производительности классической реализации и реализации на основе библиотеки GraphBLAS.

\textbf{Задачи:}
\begin{itemize}
    \item Реализовать классический BFS на основе очереди с использованием CSR-представления графа.
    \item Реализовать BFS на основе библиотеки SuiteSparse:GraphBLAS с использованием умножения вектора фронта на матрицу смежности.
    \item Реализовать загрузчик графов из формата Matrix Market (.mtx).
    \item Реализовать таймер для высокоточного измерения времени выполнения (точность до микросекунд).
    \item Покрыть реализации тестами с использованием графов известной структуры.
    \item Провести экспериментальное исследование производительности на графах из коллекции SuiteSparse Matrix Collection.
\end{itemize}

\section{Описание предлагаемого решения}

\subsection{Реализация BFS}

В рамках работы реализованы две версии BFS: с построением дерева обхода (parent BFS) и обход из нескольких источников (multi-source BFS).

\subsubsection{Классическая версия}

\begin{itemize}
    \item CSR-представление графа (массивы \texttt{row\_ptr}, \texttt{col\_idx})
    \item Обход с помощью очереди
    \item Сложность O(V + E), память O(V + E)
\end{itemize}

\subsubsection{GraphBLAS версия}

\begin{itemize}
    \item Булева разреженная матрица \texttt{GrB\_Matrix}
    \item Обход через умножение вектора фронта на матрицу смежности
    \item Полукольцо \texttt{GxB\_LOR\_LAND\_BOOL}
\end{itemize}

Реализация выполнена на языке C, сборка через CMake.

\subsection{Измерение производительности}

Для измерения времени выполнения BFS использовался таймер на основе системного вызова \verb|gettimeofday|, обеспечивающий точность до микросекунд. Принцип работы таймера:

\begin{enumerate}
    \item Запоминается время непосредственно перед вызовом функции BFS.
    \item Выполняется BFS.
    \item Запоминается время после завершения BFS.
    \item Вычисляется разница в микросекундах.
\end{enumerate}

Для повышения точности и исключения случайных флуктуаций системы каждый эксперимент проводился 10 раз. По полученным значениям вычислялись среднее арифметическое и стандартное отклонение.

\subsection{Выбранные для тестов графы}

Для проведения экспериментов были выбраны три графа из коллекции SuiteSparse Matrix Collection:

\begin{itemize}
    \item \textbf{delaunay\_n16} --- триангуляция Делоне (65 536 вершин, 196 575 рёбер). Характеризуется равномерной степенью вершин (~6) и средним диаметром (119).
    
    \item \textbf{com-Amazon} --- социальная сеть ко-покупок Amazon (334 863 вершин, 925 872 рёбер). Имеет структуру сообществ, диаметр 44.
    
    \item \textbf{kron\_g500-logn16} --- синтетический безмасштабный граф, сгенерированный по модели RMAT (65 536 вершин, 2 456 398 рёбер). Степень вершин распределена по степенному закону, диаметр всего 6.
\end{itemize}

\subsection{Интерфейс командной строки}

Программа \verb|bfs_analysis| содержит точку входа с аргументами командной строки:

\begin{verbatim}
./bfs_analysis <graph_file.mtx> [source_node]
\end{verbatim}

\begin{itemize}
    \item \verb|graph_file.mtx| --- путь к файлу в формате Matrix Market (.mtx), содержащему граф.
    \item \verb|source_node| --- (необязательный) номер стартовой вершины в 0-индексации. По умолчанию используется 0.
\end{itemize}

Примеры запуска:
\begin{verbatim}
./bfs_analysis ../graphs/delaunay_n16.mtx 27894
./bfs_analysis ../graphs/com-Amazon.mtx 334485
./bfs_analysis ../graphs/kron_g500-logn16.mtx 36114
\end{verbatim}

Результаты выполнения сохраняются в файлы:
\begin{itemize}
    \item \verb|classic_bfs_result.txt| --- результаты классического BFS (массив родителей и уровней).
    \item \verb|graphblas_bfs_result.txt| --- результаты BFS на основе GraphBLAS.
    \item \verb|bfs_comparison.txt| --- сводка времени выполнения и сравнение производительности.
\end{itemize}

\section{Функциональное тестирование}

Тесты в файлах \verb|test_classic_bfs.c| и \verb|test_graphblas_bfs.c| проверяют корректность обеих реализаций BFS на графах известной структуры:

\textbf{Классическая реализация (test\_classic\_bfs.c):}
\begin{itemize}
    \item \textbf{Линейный граф (0-1-2-3)}: проверяется правильность построения BFS-дерева (parent и level для каждого узла) и количество посещённых узлов.
    \item \textbf{Несвязный граф (0-1 и 2-3)}: проверяется корректная обработка недостижимых вершин (parent = -1, level = -1).
    \item \textbf{Граф из одной вершины}: проверяется работа на граничных условиях (один узел, нет рёбер).
\end{itemize}

\textbf{GraphBLAS реализация (test\_graphblas\_bfs.c):}
\begin{itemize}
    \item \textbf{Линейный граф (0-1-2-3)}: проверяется корректность работы GraphBLAS BFS, включая parent и level для источника и первого уровня.
    \item \textbf{Граф из одной вершины}: проверяется работа на граничных условиях (один узел, нет рёбер).
\end{itemize}

Все тесты успешно проходят (100\%), что подтверждает корректность обеих реализаций. Запуск тестов осуществляется через систему \verb|ctest|:

\begin{verbatim}
ctest
100% tests passed, 0 tests failed out of 2
\end{verbatim}

\section{Экспериментальное исследование производительности}

Эксперименты проводились на системе со следующими характеристиками: 
\textit{процессор} Intel Core i5-13420H, \textit{ОЗУ} 16~ГБ, \textit{ОС} Ubuntu 24.04.4.

Программная реализация выполнена на языке C. Исследование производительности сравнивает две реализации BFS:

\begin{itemize}
    \item \textbf{classic}~--- классическая реализация на основе очереди с использованием CSR-представления графа.
    \item \textbf{graphblas}~--- реализация на основе библиотеки SuiteSparse:GraphBLAS с использованием умножения вектора фронта на матрицу смежности.
\end{itemize}

Тесты запускаются на трёх графах из коллекции SuiteSparse Matrix Collection:
\begin{itemize}
    \item \textbf{delaunay\_n16} (65 536 вершин, 196 575 рёбер)~--- структурированный граф с равномерной степенью вершин.
    \item \textbf{com-Amazon} (334 863 вершин, 925 872 рёбер)~--- реальный граф социальной сети.
    \item \textbf{kron\_g500-logn16} (65 536 вершин, 2 456 398 рёбер)~--- синтетический безмасштабный граф.
\end{itemize}

Для каждого графа было выполнено по 10 запусков BFS от вершины с максимальной степенью. Измерялось время выполнения в микросекундах. Результаты представлены в таблице~\ref{tab:results}.

\begin{table}[ht]
    \centering
    \begin{tabular}{|p{3cm}|p{2.5cm}|p{2.5cm}|p{3.5cm}|}
        \hline
        \textbf{Граф} & \textbf{Classic BFS (мкс)} & \textbf{GraphBLAS BFS (мкс)} & \textbf{Результат} \\
        \hline
        delaunay\_n16 & $490 \pm 80$ & $178 \pm 30$ & GraphBLAS в 2.75x быстрее \\
        \hline
        com-Amazon & $2\,900 \pm 150$ & $3\,100 \pm 350$ & Classic на 7\% быстрее \\
        \hline
        kron\_g500-logn16 & $3\,350 \pm 150$ & $4\,300 \pm 300$ & Classic на 28\% быстрее \\
        \hline
    \end{tabular}
    \caption{Результаты сравнительного тестирования производительности BFS. Указаны средние значения и стандартные отклонения по 10 запускам.}
    \label{tab:results}
\end{table}

Как видно из таблицы, на структурированном графе \textbf{delaunay\_n16} реализация на GraphBLAS оказалась быстрее классической в 2.75 раза. Это объясняется равномерной структурой графа и широкими фронтами BFS, которые эффективно обрабатываются матричными операциями.

На реальном графе \textbf{com-Amazon} оба подхода показали сопоставимые результаты с небольшим преимуществом классической реализации (7\%). Это связано со структурой сообществ, которая не даёт явного преимущества ни одному из подходов.

На синтетическом безмасштабном графе \textbf{kron\_g500-logn16} классическая реализация оказалась быстрее на 28\%. Причина~--- очень маленький диаметр графа (всего 6 уровней), при котором накладные расходы на матричные операции в GraphBLAS не окупаются.

Все тесты пройдены успешно (100\%), результаты обеих реализаций согласованы (проверка на совпадение множества достижимых узлов дала положительный результат).

Таким образом, экспериментально подтверждена гипотеза о том, что выбор оптимальной реализации BFS зависит от топологии графа: GraphBLAS эффективен на структурированных графах с широкими фронтами, тогда как классическая реализация показывает лучшие результаты на графах с малым диаметром.

\section{Заключение}

В ходе выполнения работы получены следующие результаты:

\begin{itemize}
    \item Реализованы две версии BFS (классическая и GraphBLAS)
    \item Проведено тестирование на трёх графах различной топологии
    \item Выявлена зависимость производительности от структуры графа
\end{itemize}

Репозиторий с исходным кодом: \url{https://github.com/MariaPentegova/graph_blas_bfs-2}

\end{document}
